\section*{Capítulo \ref{ch:b1:erosion-transport} --- Erosión y transporte}

\begin{solution}{exo:b1:erosion-transport:1}
\qty{0.4}{g.L^{-1}} son \qty{0.4}{kg.m^{-3}}, así que el flujo es $500 \times 0.4 =
\qty{200}{kg.s^{-1}}$. Un año de \num{365.25}~días dura
\qty{3.156e7}{s}, así que el flujo anual es \qty{6.31e9}{kg}, o
\qty{6.31}{Mt}.
\end{solution}

\begin{solution}{exo:b1:erosion-transport:2}
$w_s = R g d^2/(18\nu) = 1.6548 \times 9.81 \times (\num{30e-6})^2/(18 \times
\num{1.0034e-6}) = \qty{8.09e-4}{m.s^{-1}}$, unos \qty{0.8}{mm} por segundo. Su
número de Reynolds del grano es $w_s d/\nu = \num{8.09e-4} \times
\num{30e-6}/\num{1.0034e-6} = 0.024$, muy por debajo de 1: vale la ley de Stokes. Se hunde
\qty{2}{m} en $2/\num{8.09e-4} = \qty{2472}{s}$, unos
\qty{41}{min}.
\end{solution}

\begin{solution}{exo:b1:erosion-transport:3}
$\tau_b = \rho g h S = 998.2 \times 9.81 \times 0.5 \times \num{4e-4} =
\qty{1.96}{Pa}$. Con $(\rho_s - \rho) g = 1651.8 \times 9.81 =
\qty{16204}{N.m^{-3}}$: para \qty{1}{mm}, $\theta = 1.96/(16204 \times 0.001) =
0.121$, casi cuatro veces el umbral de \num{0.0314}, así que la arena se mueve; para
\qty{10}{mm}, $\theta = 0.0121$, una quinta parte del umbral de \num{0.0556}, así que la
grava permanece. El arroyo transporta su arena sobre un lecho de grava que no se mueve.
\end{solution}

\begin{solution}{exo:b1:erosion-transport:4}
(a)~$12.6/3.9 = 3.2$: los ríos transportan en suspensión unas tres veces más que en disolución. (b)~Grandes embalses: $0.20 \times 14 = \qty{2.8}{Gt}$ al
año; todos ellos: $0.26 \times 14 = \qty{3.64}{Gt}$ al año. (c)~El flujo disuelto es de \qty{3.9e15}{g} al año a \qty{120}{g.m^{-3}}, así que el agua
que lo lleva es $\num{3.9e15}/120 = \qty{3.25e13}{m^3}$, unos \qty{32500}{km^3} al año. La \omterm{def:b1:erosion-transport:load}{carga en suspensión} en esa agua es $\qty{1.26e16}{g}/\qty{3.25e16}{L} =
\qty{0.39}{g.L^{-1}}$.
\end{solution}

\begin{solution}{exo:b1:erosion-transport:5}
(a)~La sección tiene el área $10 \times 2 = \qty{20}{m^2}$ y el perímetro mojado $10 + 2 \times 2 = \qty{14}{m}$, así que $R_h = 20/14 = \qty{1.43}{m}$.
(b)~$\tau_b = 998.2 \times 9.81 \times 1.43 \times 0.001 = \qty{14.0}{Pa}$,
frente a $\qty{19.6}{Pa}$ para un canal muy ancho de \qty{2}{m} de profundidad: el cociente es
$R_h/h = 0.714$, pues las orillas toman el resto del rozamiento. (c)~Para un canal rectangular de
anchura $b$, $R_h = bh/(b + 2h) = h/(1 + 2h/b)$, que tiende a $h$ cuando $b$ crece.
\end{solution}

\begin{solution}{exo:b1:erosion-transport:6}
(a)~$d_1 = \left(18\nu^2/(Rg)\right)^{1/3} = (18/(Rg))^{1/3}\,\nu^{2/3}$, y $R$
y $g$ no dependen de la viscosidad. (b)~$\nu = \num{1.3059e-3}/999.7 =
\qty{1.306e-6}{m^2.s^{-1}}$ a \qty{10}{\celsius}, frente a
\qty{1.0034e-6}{m^2.s^{-1}} a \qty{20}{\celsius}: el cociente es \num{1.302},
y $1.302^{2/3} = 1.192$, así que $d_1 = 0.104 \times 1.192 = \qty{0.124}{mm}$.
(c)~En el intervalo de Stokes $w_s$ es inversamente proporcional a $\nu$: el grano
se hunde $1.302$ veces más despacio, unos \qty{30}{\percent}. El agua fría mantiene
el sedimento fino en suspensión durante más tiempo: un río en invierno transporta su limo
algo más lejos para el mismo caudal.
\end{solution}

\begin{solution}{exo:b1:erosion-transport:7}
Con $R g/(18\nu) = \qty{8.99e5}{m^{-1}.s^{-1}}$ y la ley combinada:
\begin{center}
\begin{tabular}{@{}lccc@{}}
\hline
$d$ & Stokes (\si{m.s^{-1}}) & combinada (\si{m.s^{-1}}) & cociente\\
\hline
\qty{0.1}{mm} & \num{8.99e-3} & \num{7.53e-3} & \num{1.19}\\
\qty{0.5}{mm} & \num{0.225} & \num{0.0711} & \num{3.16}\\
\qty{2}{mm} & \num{3.60} & \num{0.197} & \num{18.3}\\
\hline
\end{tabular}
\end{center}
El cociente crece con el tamaño, porque crece el número de Reynolds del grano: la ley de Stokes, válida solo para números de Reynolds pequeños,
aumenta como $d^2$, mientras que el arrastre de la estela frena los granos gruesos hasta una velocidad que crece solo como $\sqrt d$.
\end{solution}

\begin{solution}{exo:b1:erosion-transport:8}
(a)~$(Rg/\nu^2)^{1/3} = \qty{25263}{m^{-1}}$. A \qty{0.2}{mm}, $D_* = 5.05$
y $\theta_c = 0.3/7.06 + 0.055(1 - e^{-0.101}) = 0.0425 + 0.0053 =
0.0478$; $\tau_c = 0.0478 \times 16204 \times \num{2e-4} = \qty{0.155}{Pa}$.
A \qty{20}{mm}, $D_* = 505$ y $\theta_c = 0.3/607 + 0.055(1 - e^{-10.1}) =
0.0005 + 0.0550 = 0.0555$; $\tau_c = 0.0555 \times 16204 \times 0.02 =
\qty{18.0}{Pa}$. (b)~$h = \tau_c/(\rho g S)$ con $\rho g S =
\qty{19.58}{Pa.m^{-1}}$: \qty{7.9}{mm} de agua mueven la arena fina, \qty{0.92}{m}
la grava.
\end{solution}

\begin{solution}{exo:b1:erosion-transport:9}
La cuenca pierde \qty{0.8}{Mt} = \qty{8e5}{t} al año, o $\num{8e5}/2000 =
\qty{400}{t.km^{-2}.yr^{-1}}$. Según \cref{prop:b1:erosion-transport:denudation},
$D = \qty{8e8}{kg}/(2650 \times \qty{2e9}{m^2}) = \qty{1.51e-4}{m}$ al año, o
\qty{151}{m/Myr}. Eso es unas tres veces la mediana de las cuencas medidas
con berilio-10, \qty{54}{m/Myr}, aunque por debajo de su media de
\qty{218}{m/Myr}: una cuenca bastante rápida, como lo son las de relieve empinado o de región tectónicamente activa.
\end{solution}

\begin{solution}{exo:b1:erosion-transport:10}
(a)~$u_* = \sqrt{g h S} = \sqrt{9.81 \times 3 \times \num{5e-4}} =
\qty{0.1213}{m.s^{-1}}$ y $\tau_b = \rho u_*^2 = \qty{14.7}{Pa}$.
(b)~La suspensión necesita $w_s < u_* = \qty{0.1213}{m.s^{-1}}$. La ley combinada
da \qty{0.126}{m.s^{-1}} a \qty{1}{mm}; resolviendo
$w_s(d) = 0.1213$ por tanteo resulta $d = \qty{0.94}{mm}$. (c)~Para la grava
$\theta_c \approx 0.0555$, así que $\tau_c = \qty{14.7}{Pa}$ para $d \approx
14.7/(0.0555 \times 16204) = \qty{16.3}{mm}$; a ese tamaño la curva de umbral da
$\theta_c = 0.0555$, en efecto. (d)~Los granos más gruesos que unos \qty{16}{mm} permanecen en el fondo y forman un suelo estable de grava; la grava más fina y la arena
más gruesa que cerca de un milímetro viajan como \omterm{def:b1:erosion-transport:load}{carga de fondo}; la arena fina, el limo y la arcilla viajan en suspensión.
\end{solution}

\begin{solution}{exo:b1:erosion-transport:11}
(a)~$0.30 \times \qty{1000}{m} = \qty{300}{m}$ de roca, promediada sobre la región. (b)~El rebote es $300 \times 2700/3300 = \qty{245}{m}$. (c)~La superficie media
baja solo $300 - 245 = \qty{55}{m}$. (d)~Las cumbres, que no fueron erosionadas, son elevadas con toda la región: ahora están
\qty{245}{m} más altas que antes. Cortar valles puede elevar los picos de una cordillera.
\end{solution}

\begin{solution}{exo:b1:erosion-transport:12}
(a)~$0.9 \times \qty{5e8}{kg}/\qty{1300}{kg.m^{-3}} = \qty{346000}{m^3}$ al año,
o \qty{0.69}{\percent} del embalse. (b)~La mitad del embalse,
\qty{25e6}{m^3}, se llena en $\num{25e6}/\num{346000} = \num{72}\text{ años}$ a este ritmo. (c)~A medida que el embalse se llena, su volumen, y con él el tiempo
de residencia del agua, disminuye, y pasan granos más finos: su eficacia de retención baja.
Mientras tanto, los depósitos más antiguos quedan enterrados y se compactan, expulsando agua, de modo que su densidad aparente sube y cada tonelada ocupa menos sitio.
Ambos efectos frenan la pérdida de volumen a medida que el embalse envejece.
\end{solution}

\begin{solution}{pb:b1:erosion-transport:1}
\textbf{1.} Un año de agua son $1.2 \times \num{3.156e7} = \qty{3.79e7}{m^3}$;
repartidos sobre \qty{8e7}{m^2}, una capa de \qty{0.473}{m}: una escorrentía de
\qty{473}{mm} al año.

\textbf{2.} $1.2 \times \qty{0.150}{kg.m^{-3}} = \qty{0.18}{kg.s^{-1}}$, o
\qty{5680}{t} al año.

\textbf{3.} $1.2 \times 0.080 = \qty{0.096}{kg.s^{-1}}$, o \qty{3030}{t} al
año.

\textbf{4.} La \omterm{def:b1:erosion-transport:load}{carga de fondo} es de \qty{568}{t} al año, y la total $5680 + 3030 +
568 = \qty{9278}{t}$ al año.

\textbf{5.} La concentración del sedimento en suspensión es alta durante los pocos días
de crecida, cuando el caudal también es alto, y baja el resto del tiempo. Un promedio temporal da a cada día el mismo peso, así que lo domina
la multitud de días de caudal bajo y concentración baja. Según
\cref{prop:b1:erosion-transport:flux}, la masa anual es $\int Q\,C\,\mathrm dt$; la media ponderada por el flujo es esta integral dividida por
$\int Q\,\mathrm dt$, así que es la media ponderada por el flujo la que debe multiplicarse por el caudal medio.

\textbf{6.} $3030/9278 = \num{0.33}$, un tercio. Para los ríos del mundo,
$3.9/(12.6 + 3.9) = \num{0.24}$, y menos si se contara la \omterm{def:b1:erosion-transport:load}{carga de fondo}.
La parte de la cuenca es mayor: colinas bajas, donde la erosión mecánica lenta deja tiempo a la roca para disolverse, y calizas y creta, que se disuelven
por completo.

\textbf{7.} $9278/80 = \qty{116}{t.km^{-2}.yr^{-1}}$.

\textbf{8.} $D = \qty{9.278e6}{kg}/(2700 \times \qty{8e7}{m^2}) =
\qty{4.30e-5}{m}$ al año, \qty{43.0}{m/Myr}.

\textbf{9.} $43.0/54 = 0.80$: un poco por debajo de la mediana de las cuencas medidas, y cinco veces por debajo de su media. Una cuenca moderadamente
lenta, como cabe esperar para colinas bajas en una región tranquila.

\textbf{10.} El aforo promedia los pocos años de medición, y depende de las crecidas que ocurrieron en ellos. El berilio-10 promedia el tiempo que se tarda en
erosionar los pocos metros en los que se produce: a \qty{43}{m/Myr}, cada metro tarda \num{23000} años, así que el promedio abarca decenas de miles de años. Pueden
diferir por las grandes crecidas raras, ausentes en un registro corto; porque la agricultura y la tala han aumentado la erosión en los últimos siglos; y porque el
sedimento se almacena en la llanura de inundación, y de allí se recupera.

\textbf{11.} Parte de la \omterm{def:b1:erosion-transport:load}{carga disuelta} no procede de la roca: buena parte del bicarbonato lleva carbono tomado del aire o del suelo, y la lluvia trae sal marina.
Contarlos acreditaría a la cuenca una roca que no ha perdido.

\textbf{12.} $\num{1.5e6}/1.2 = \qty{1.25e6}{s}$, \num{14.5} días.

\textbf{13.} La ley combinada da \qty{7.53e-3}{m.s^{-1}}: el grano cae
\qty{8}{m} en \qty{1062}{s}, unos \qty{18}{min}. Queda atrapado.

\textbf{14.} $w_s = \num{8.99e5} \times (\num{2e-5})^2 =
\qty{3.60e-4}{m.s^{-1}}$: \qty{8}{m} en \qty{22200}{s}, unas \num{6.2}~horas.
Atrapado.

\textbf{15.} $w_s = \num{8.99e5} \times (\num{1e-6})^2 =
\qty{8.99e-7}{m.s^{-1}}$: \qty{8}{m} tardan \qty{8.9e6}{s}, \num{103}~días, siete
veces el tiempo de residencia. La mayor parte pasa por la presa con el agua.

\textbf{16.} El grano debe caer \qty{8}{m} en \qty{1.25e6}{s}, es decir, a $\qty{6.4e-6}{m.s^{-1}}$ como mínimo. Según la ley de Stokes $d = \sqrt{18\nu
w_s/(Rg)} = \sqrt{\num{6.4e-6}/\num{8.99e5}} = \qty{2.7}{\micro m}$: una arcilla fina. Los granos más gruesos quedan atrapados donde entren.

\textbf{17.} $568 + 0.7 \times 5680 = \qty{4544}{t}$ al año, que llenan
$\num{4.544e6}/1300 = \qty{3496}{m^3}$ al año, el \qty{0.233}{\percent} del
embalse. La mitad, \qty{7.5e5}{m^3}, se llena en unos \num{215} años.

\textbf{18.} $100/43.0 = \qty{2.33}{Myr}$.

\textbf{19.} $1 - 2700/3300 = 0.18$: la superficie baja solo el
\qty{18}{\percent} del espesor retirado.

\textbf{20.} $100/0.1818 = \qty{550}{m}$ de roca.

\textbf{21.} Toda la región se eleva $550 \times 2700/3300 = \qty{450}{m}$ por rebote. Las cumbres apenas erosionadas son elevadas con ella, hasta
\qty{450}{m}, mientras los valles se profundizan: el relieve crece antes de disminuir.

\textbf{22.} La \omterm{def:b1:erosion-transport:denudation}{tasa de denudación} depende sobre todo de las pendientes. A medida que las colinas se rebajan y sus pendientes se suavizan, la tasa
disminuirá, salvo que el levantamiento o un cambio de clima las vuelva a empinar: el tiempo calculado a la tasa de hoy es una cota inferior.

\textbf{23.} $550/42.95 = 12.80$. A la tasa de hoy y con el rebote isostático,
bajar las Varden Hills \qty{100}{m} tarda \textbf{12.8}~millones de años.
\end{solution}
